\chapter{السنن الناسخة والمنسوخة وتخصيص الأحكام}

\section{نسخ النهي عن ادخار لحوم الأضاحي}
الحمد لله والصلاة والسلام على رسول الله، أما بعد؛ فقد استكمل الإمام الشافعي رحمه الله تعالى الحديث عن موقع السنة من القرآن، وبدأ يفصل في "السنن الناسخة والمنسوخة". فكما أن القرآن ينسخ بعضه بعضاً، فكذلك سنة رسول الله صلى الله عليه وسلم ينسخ بعضها بعضاً.

ومن ذلك نهي النبي صلى الله عليه وسلم عن ادخار لحوم الأضاحي بعد ثلاث. فقد روى ابن عمر وعلي بن أبي طالب رضي الله عنهما أن رسول الله صلى الله عليه وسلم نهى عن أكل لحوم الضحايا بعد ثلاث. وقد حفظا المنسوخ وحدثا به، ولم تبلغهما الرخصة الناسخة.
بينما حفظت أم المؤمنين عائشة رضي الله عنها المنسوخ والناسخ وسبب التشريع فقالت: «دَفَّ ناسٌ من أهل البادية حضرة الأضحى في زمان النبي صلى الله عليه وسلم (والدافة: هم القوم يسيرون جماعة ونزلوا ضيوفاً)، فقال النبي صلى الله عليه وسلم: ادخروا لثلاث وتصدقوا بما بقي. فلما كان بعد ذلك قيل يا رسول الله لقد كان الناس ينتفعون بضحاياهم يُجملون منها الودك (الشحم) فقال: إنما نهيتكم من أجل الدافة التي دفت حضرة الأضحى، فكلوا وتصدقوا وادخروا».

فحديث عائشة رضي الله عنها من أبين ما يوجد في الناسخ والمنسوخ، وهو يوضح أن بعض الرواة قد يحفظ المنسوخ ولا يبلغه الناسخ، ولا تضيع السنة بتمامها عن مجموع الأمة. 
وقد رجّح المحققون (كالعلامة أحمد شاكر) أن نهي النبي صلى الله عليه وسلم هنا لم يكن تشريعاً عاماً منسوخاً، بل كان تصرفاً منه على سبيل نظر الإمام والحاكم لمصلحة الرعية لوجود المجاعة (الدافة)، فلما زالت العلة زال النهي.

\section{صلاة الخوف ونسخ تأخير الصلاة عن وقتها}
من أمثلة الناسخ والمنسوخ في السنة: تأخير الصلاة عن وقتها في شدة الخوف.
فقد روى أبو سعيد الخدري رضي الله عنه أن المشركين شغلوا النبي صلى الله عليه وسلم يوم الخندق عن صلاة الظهر والعصر والمغرب حتى دخل الليل، فأمر بلالاً فأقام وصلاها تباعاً. وكان هذا قبل نزول شرعية (صلاة الخوف).

فلما نزل قوله تعالى: ﴿وَإِذَا كُنتَ فِيهِمْ فَأَقَمْتَ لَهُمُ الصَّلَاةَ...﴾ شرع النبي صلى الله عليه وسلم صلاة الخوف بكيفيات متعددة (كصلاة ذات الرقاع وعسفان)، فصارت الصلاة تؤدى في وقتها على أي هيئة أمكنت (رجالاً أو ركباناً، مستقبلي القبلة أو غير مستقبليها في حال المسايفة). وبذلك نُسخت رخصة تأخير الصلاة عن وقتها لعذر الخوف، ولا يجوز تأخيرها بحال بل تُصلى في وقتها قدر الاستطاعة.

\section{نسخ الحبس في الزنا وتفصيل إقامة الحدود}
كان عقاب الزاني في أول الإسلام الحبس والأذى، كما قال تعالى: ﴿وَاللَّاتِي يَأْتِينَ الْفَاحِشَةَ مِن نِّسَائِكُمْ فَاسْتَشْهِدُوا عَلَيْهِنَّ أَرْبَعَةً مِّنكُمْ ۖ فَإِن شَهِدُوا فَأَمْسِكُوهُنَّ فِي الْبُيُوتِ حَتَّىٰ يَتَوَفَّاهُنَّ الْمَوْتُ أَوْ يَجْعَلَ اللَّهُ لَهُنَّ سَبِيلًا﴾. 
ثم نسخ الله هذا الحبس بآية سورة النور: ﴿الزَّانِيَةُ وَالزَّانِي فَاجْلِدُوا كُلَّ وَاحِدٍ مِّنْهُمَا مِائَةَ جَلْدَةٍ﴾. 
وبينت السنة المطهرة في حديث عبادة بن الصامت هذا "السبيل"، فقال صلى الله عليه وسلم: «خذوا عني، قد جعل الله لهن سبيلاً؛ البكر بالبكر جلد مائة وتغريب عام، والثيب بالثيب جلد مائة والرجم». ثم نُسخ الجلد عن الثيب واكتفي بالرجم فقط كما فعل النبي صلى الله عليه وسلم في رجم ماعز والغامدية.

\subsection{تنصيف الحد للإماء}
قال تعالى في الإماء: ﴿فَإِذَا أُحْصِنَّ فَإِنْ أَتَيْنَ بِفَاحِشَةٍ فَعَلَيْهِنَّ نِصْفُ مَا عَلَى الْمُحْصَنَاتِ مِنَ الْعَذَابِ﴾. 
وقد بين الشافعي أن التنصيف لا يكون إلا فيما يقبل العدد وهو (الجلد)، فتجلد الأمة خمسين جلدة. أما الحدود المتلفة للنفوس كـ (الرجم) فلا يمكن تنصيفها؛ لأن المرجوم قد يموت بحجر واحد أو بألف، وما لا يقبل العدد لا يقبل التنصيف، ولذلك أجمع العلماء على أن الأمة إذا زنت لا تُرجم.

\section{دقة الشريعة البالغة في الدماء والأعراض (قصة العسيف)}
لإثبات عظمة الشريعة واحتياطها في الدماء والأعراض، أورد الشافعي قصة (العسيف)، وهو الأجير الذي زنى بامرأة مستأجره، فافتداه أبوه بمائة شاة وجارية ظناً منه أن ابنه سيُرجم، فلما ارتفعوا للنبي صلى الله عليه وسلم، أبطل هذه الفدية لعدم شرعيتها، وأمر بجلد الشاب مائة وتغريبه عاماً (لأنه بكر غير محصن). 

وأما المرأة المشتكى عليها، فقد قال النبي صلى الله عليه وسلم لأُنيس الأسلمي: «واغدُ يا أُنيس إلى امرأة هذا، فإن اعترفت فارجمها». وفي هذا رد مفحم على أعداء الشريعة؛ فرغم اعتراف الشريك (الشاب)، وإقرار الزوج، وافتداء أب الشاب بالمال، لم يقم النبي صلى الله عليه وسلم الحد على المرأة فوراً! بل اشترط إما أن تعترف هي بنفسها، أو توجد البينة القطعية بـ (أربعة شهداء)، فاعترفت فرجمها، ولو أنكرت لدُرئ عنها الحد.

\subsection{خضوع الذميين لقانون العقوبات الإسلامي}
روى الشافعي عن ابن عمر رضي الله عنهما أن النبي صلى الله عليه وسلم رجم يهوديين زنيا. وهذا دليل ساطع على أن الشريعة الإسلامية هي الحاكمة والنافذة في ديار الإسلام على جميع المقيمين فيها والمواطنين، حتى من غير المسلمين (كأهل الذمة)، إذا ترافعوا إلينا أو أعلنوا منكرهم، ولا حصانة لقانون كفري يخالف شرع الله داخل دولة الإسلام.

\section{صلاة المأموم خلف الإمام القاعد (بين النسخ والإحكام)}
من أسباب الاختلاف بين العلماء في تصحيح الأحاديث دعوى النسخ في صلاة المأموم خلف الإمام إذا صلى قاعداً لعذر.
فقد ثبت في حديث أنس أن النبي صلى الله عليه وسلم سقط من فرس فجُحش شقه (أي جُرح)، فصلى قاعداً وصلى الصحابة وراءه قعوداً، وقال: «وإذا صلى جالساً فصلوا جلوساً أجمعون».
ثم ثبت في حديث عائشة في مرض موته صلى الله عليه وسلم، أنه خرج فصلى قاعداً إلى جنب أبي بكر، وكان الناس يصلون قياماً خلف أبي بكر.

وقد ذهب الإمام الشافعي إلى أن فعل النبي صلى الله عليه وسلم في مرض موته ناسخ لأمره الأول بجلوس المأمومين، وأن من أطاق القيام لا يجوز له أن يجلس خلف الإمام المريض.
بينما ذهب فقهاء آخرون (كالإمام أحمد، ورجحه المحقق أحمد شاكر) إلى عدم النسخ، لأن الأمر بالقعود جاء بأعلى صيغ التأكيد والنهي عن التشبه بالأعاجم، وأن صلاة أبي بكر بالناس قياماً في مرض الموت تحتمل أنه كان هو الإمام المُسمِع للناس، وهي حالة خاصة. والخلاف في ذلك معتبر ومبني على دقة توجيه النصوص.